\section{Reproducibility and source identities}\label{app:provenance}
\subsection{Recorded execution and failures}
Versioned records retain inputs, configurations, commands, dependencies, code revisions, logs and every outcome. E1--E5 correspond respectively to \path{recovery_v0.8}, \path{commissioning_v0.9}, \path{casper_convergence_v0.11}, \path{return_bound_v0.12}, and \path{residual_twin_v0.13} followed by \path{attack_audit_v0.14}. Their calibration policies are not interchangeable. No experiment was rerun for this article revision.

E1 consumed 229.20 s wall, 228.07 s CPU and 167.25 MiB sampled memory; E2 consumed 14.75 s, 14.30 s and 113.97 MiB. E1's initial launch failed before initialization because monitoring was blocked. E2 stopped once on a reporting-key collision after selection; the recorded repair preserved identical fits before fresh generation. E5 has no failed scientific attempt. Its final limits were 120 s wall, 90 s CPU, 1 GiB monitored memory and 128 MiB output. Worker peak memory was 99.53 MiB versus 90.17 MiB sampled; sampling is not an instantaneous hard bound. The IMU run used 2.68 MiB outputs and passed 17 synthetic tests and 3,754 saved-artifact checks. E4 used 0.565 s wall, 0.459 s CPU, 39.81 MiB memory and 11.07 MiB outputs, with 20 tests and 2,665 checks.

The final freeze SHA-256 is \path{b48a0c3927b465559a8f7b372c2d26ed7148642c97d2435de83995a0ad3be021}; its result-package manifest is \path{987d8904dda988c7ccd51d61a5e52cf8e162e8f0156286aad89c385813b785dc}. Asset receipts bind figures and tables to saved results. Identifiers establish provenance, not validity. Raw data and third-party manuscripts are excluded from the editable archive.

\subsection{Real-data identity and comparator admission}
The author-linked CASPER-1-and-2 provider-v2 nominal IMU file has SHA-256 \path{1221df9be80c229969f3f21f820837e68ad28f4ccb3949bf4013f00d6d69d62c}. Its original-paper run and clock correspondence remain unresolved.

The original voraus-AD file is linked by \url{https://github.com/vorausrobotik/voraus-ad-dataset}. It has 1,115,942,833 bytes and SHA-256 \path{c90ab1c78af52651b954d41787f7e89d750f0a128b57600b0e5ceec22621f704}. Data are CC BY-NC-SA 4.0; code is MIT. Metadata verify 2,321,690 rows, 137 columns and the 2,122 executions in Table~\ref{tab:data}. Setting 72 (PRE\_A) supplies all 948 training executions. Within-execution time resets to zero, increases strictly with increments 0.00845194--0.01211095 s, and does not establish absolute acquisition order. Maximum train/test lengths are both 1,164 rows. Process signals remain unopened.

Pinned MVT-Flow code \path{a91a86a642d23df58833b792a53de01edfd81abe} has four coupling blocks, clamp 1.2, width scale two, Adam $8\times10^{-4}$, batch 32 and 70 epochs. It computes test AUROC each epoch but saves the fixed final epoch; this alone does not prove test-based selection. Training uses $\|z\|^2-\log|J|$, while paper/scoring use $\tfrac12\|z\|^2-\log|J|$. The paper describes three convolutions; the default released block constructs two. A padding truncation is overwritten, but the received release has no overlength test execution, so no impact is demonstrated. These findings leave faithful reproduction unresolved; they do not prove invalid published results.

FactoryNet arXiv:2605.09081v5 (8 October 2026) reports ICML-associated workshop acceptance, not verified main-conference publication. The pinned data card \path{3bea3cfc10545d898604f93170108576fd46d518} reports 56,591 executions: 14,847 real and 41,744 simulated. Code revision is \path{651f066d97b704f0824687f1086eeb8f0422a8f3}. The card says Isaac Sim 5.1 data are still uploading; the available snapshot is 4.5. Imported voraus-AD, AURSAD and CNC are not independent new acquisitions. Context includes fault labels, channels vary by embodiment, and licenses differ by subset/code. No observations were downloaded or performance measured.
